% Section 1 Header
\section{Section 1 — Two Pointers: The Simplest Powerful Pattern}

\begin{frame}[c]{\empty}
    \begin{center}
        \Large\textbf{\secname}
    \end{center}
\end{frame}

% Slide 1: Introduction to the Pattern
\begin{frame}{What is the Two Pointers Pattern?}
  \begin{itemize}
    \item \textbf{The Core Idea:} Use two indices (pointers) to iterate through a data structure simultaneously.
    \item \textbf{Why it works:} It reduces $O(n^2)$ nested loops to a single $O(n)$ pass.
    \item \textbf{Common Triggers:}
      \begin{itemize}
        \item The array is \textbf{Sorted}.
        \item You need to find a pair, a triplet, or a subarray.
        \item You need to "filter" or "rearrange" an array in-place.
      \end{itemize}
  \end{itemize}
\end{frame}

% Slide 2: How it Works (Visual Logic)
\begin{frame}{Pattern Mechanics: Opposite Ends}
  When searching for a target in a \textbf{Sorted} array:
  \begin{enumerate}
    \item \textbf{Left Pointer:} Starts at index $0$.
    \item \textbf{Right Pointer:} Starts at index $n-1$.
    \item \textbf{Decision:}
      \begin{itemize}
        \item If \textit{sum > target}: Move Right pointer left ($\leftarrow$) to decrease sum.
        \item If \textit{sum < target}: Move Left pointer right ($\rightarrow$) to increase sum.
      \end{itemize}
  \end{enumerate}
  
\end{frame}

% Slide 3: Two Sum II - Kotlin Implementation
\begin{frame}[fragile]{Hero Problem: Two Sum (Sorted)}
  \begin{minted}[fontsize=\scriptsize, frame=lines]{kotlin}
fun twoSum(numbers: IntArray, target: Int): IntArray {
    var left = 0
    var right = numbers.size - 1

    while (left < right) {
        val sum = numbers[left] + numbers[right]
        when {
            sum == target -> return intArrayOf(left + 1, right + 1)
            sum < target -> left++
            else -> right--
        }
    }
    return intArrayOf()
}
  \end{minted}
  \textbf{Time Complexity:} $O(n)$ \\
  \textbf{Space Complexity:} $O(1)$
\end{frame}

% Slide 4: Key Takeaways
\begin{frame}{Two Pointers – Key Takeaways}
  \begin{itemize}
    \item \emoji{check-mark-button} \textbf{Space Efficiency:} Usually allows for $O(1)$ extra space.
    \item \emoji{check-mark-button} \textbf{In-place operations:} Perfect for "Move Zeroes" or "Remove Duplicates."
    \item \emoji{check-mark-button} \textbf{Pointer Variants:}
      \begin{itemize}
        \item \textbf{Opposite Ends:} (Left, Right) $\to$ Searching/Sorting.
        \item \textbf{Same Direction:} (Slow, Fast) $\to$ Removing elements.
      \end{itemize}
  \end{itemize}
\end{frame}